%! Elimination Matrices and Inverse Matrices — Unit 2, Lesson 2.2 — Homework (Answer Key)
\documentclass[10pt]{article}
\usepackage{linalg-article}
\usepackage{linalg-key}   % pulls in -boxes; do NOT also load -boxes
\newcommand{\TallMath}[1]{$\displaystyle #1\rule[-1.4em]{0pt}{3.2em}$}
\newcommand{\xx}{\mathbf{x}}
\newcommand{\bb}{\mathbf{b}}

\begin{document}

\pageheader{Unit 2, Lesson 2.2}{Homework --- Answer Key}
\namedateperiod

\vspace{0.12in}

\begin{notesbox}{Practice}
Show your steps. Where you find an inverse, \textbf{check} it by multiplying $A^{-1}A=I$.

\begin{enumerate}[leftmargin=1.9em, itemsep=11pt]
  \item \textbf{Step matrices.} For each $A$, name the pivot and multiplier $\ell$, then write the
        elimination matrix $E_{21}$ and its undo-matrix $E_{21}^{-1}$.
    \begin{enumerate}[label=(\alph*), leftmargin=1.8em, itemsep=4pt]
      \item \ans{pivot $1$, $\ell=3$; $E_{21}=\begin{bsmallmatrix} 1&0\\-3&1 \end{bsmallmatrix}$, $E_{21}^{-1}=\begin{bsmallmatrix} 1&0\\3&1 \end{bsmallmatrix}$} \hfill
      \item \ans{pivot $2$, $\ell=3$; $E_{21}=\begin{bsmallmatrix} 1&0\\-3&1 \end{bsmallmatrix}$, $E_{21}^{-1}=\begin{bsmallmatrix} 1&0\\3&1 \end{bsmallmatrix}$}
    \end{enumerate}
  \item \textbf{Do the step.} For the two matrices in Problem~1, multiply $E_{21}A$ to reach
        upper-triangular $U$.
        \ansline{(a) $E_{21}A=\begin{bsmallmatrix} 1&4\\0&-7 \end{bsmallmatrix}$ \ ($-3(1,4)+(3,5)=(0,-7)$).}
        \ansline{(b) $E_{21}A=\begin{bsmallmatrix} 2&1\\0&2 \end{bsmallmatrix}$ \ ($-3(2,1)+(6,5)=(0,2)$).}
  \item \textbf{Inverses by formula.} Find $A^{-1}=\dfrac{1}{ad-bc}\begin{bsmallmatrix} d&-b\\-c&a \end{bsmallmatrix}$
        for each, and check $A^{-1}A=I$.
    \begin{enumerate}[label=(\alph*), leftmargin=1.8em, itemsep=4pt]
      \item \ans{$ad-bc=1$; $A^{-1}=\begin{bsmallmatrix} 3&-1\\-5&2 \end{bsmallmatrix}$} \hfill
      \item \ans{$ad-bc=2$; $A^{-1}=\tfrac12\begin{bsmallmatrix} 4&-3\\-2&2 \end{bsmallmatrix}=\begin{bsmallmatrix} 2&-\frac32\\-1&1 \end{bsmallmatrix}$}
    \end{enumerate}
    {\footnotesize Checks: (a) $\begin{bsmallmatrix} 3&-1\\-5&2 \end{bsmallmatrix}\begin{bsmallmatrix} 2&1\\5&3 \end{bsmallmatrix}=I$;\ \
    (b) $\begin{bsmallmatrix} 2&-3/2\\-1&1 \end{bsmallmatrix}\begin{bsmallmatrix} 2&3\\2&4 \end{bsmallmatrix}=I$.}
  \item \textbf{Juice bar --- one machine, two batches.} Each liter of \textbf{Mix P} carries $1$ unit
        sugar and $2$ units vitamin; each liter of \textbf{Mix Q} carries $2$ sugar and $5$ vitamin.
        With $p$ liters of P and $q$ of Q the totals are $A\begin{bsmallmatrix} p\\q \end{bsmallmatrix}=\bb$,
        $A=\begin{bmatrix} 1 & 2 \\ 2 & 5 \end{bmatrix}$ (sugar row, vitamin row).
    \begin{enumerate}[label=(\alph*), leftmargin=1.8em, itemsep=6pt]
      \item Build $A^{-1}$ once. \ans{$ad-bc=1$, so $A^{-1}=\begin{bsmallmatrix} 5&-2\\-2&1 \end{bsmallmatrix}$.}
      \item \textbf{Batch 1:} sugar $8$, vitamin $19$. Solve $\xx=A^{-1}\bb$ for $p,q$.
            \ansline{$\begin{bsmallmatrix} 5&-2\\-2&1 \end{bsmallmatrix}\begin{bsmallmatrix} 8\\19 \end{bsmallmatrix}=\begin{bsmallmatrix} 40-38\\-16+19 \end{bsmallmatrix}=\begin{bsmallmatrix} 2\\3 \end{bsmallmatrix}$: $p=2$, $q=3$.}
      \item \textbf{Batch 2:} sugar $9$, vitamin $21$. Same $A^{-1}$ --- solve for $p,q$.
            \ansline{$\begin{bsmallmatrix} 5&-2\\-2&1 \end{bsmallmatrix}\begin{bsmallmatrix} 9\\21 \end{bsmallmatrix}=\begin{bsmallmatrix} 45-42\\-18+21 \end{bsmallmatrix}=\begin{bsmallmatrix} 3\\3 \end{bsmallmatrix}$: $p=3$, $q=3$.}
      \item \emph{Interpret:} give the liters for each batch, and say why one inverse handled both.
            \ansline{Batch 1: $2$ L P, $3$ L Q; Batch 2: $3$ L P, $3$ L Q. The mixes (so $A$) never
            change, so the one machine $A^{-1}$ solves any batch in a single multiplication.}
    \end{enumerate}
  \item \textbf{Gauss--Jordan.} Find $A^{-1}$ for $A=\begin{bmatrix} 1 & 2 \\ 4 & 9 \end{bmatrix}$ by
        reducing $[\,A\mid I\,]$ down then up to $[\,I\mid A^{-1}\,]$.
        \[
          \left[\begin{array}{cc|cc} 1 & 2 & 1 & 0 \\ 4 & 9 & 0 & 1 \end{array}\right]
        \]
        \ansline{row$_2-4\,$row$_1$: $\left[\begin{smallmatrix} 1&2&1&0\\0&1&-4&1 \end{smallmatrix}\right]$;
        then row$_1-2\,$row$_2$: $\left[\begin{smallmatrix} 1&0&9&-2\\0&1&-4&1 \end{smallmatrix}\right]$.
        So $A^{-1}=\begin{bsmallmatrix} 9&-2\\-4&1 \end{bsmallmatrix}$.}
  \item \textbf{No inverse.} Show that $A=\begin{bmatrix} 3 & 6 \\ 2 & 4 \end{bmatrix}$ is
        \textbf{singular}: compute $ad-bc$, run one elimination step, and explain why no inverse
        exists.
        \ansline{$ad-bc=12-12=0$, so the $\tfrac{1}{ad-bc}$ formula fails. Elimination
        (row$_2-\tfrac23\,$row$_1$) gives $(0,0)$ --- a zero pivot. $A$ is singular: no inverse
        (its rows/columns are proportional).}
\end{enumerate}
\end{notesbox}

\begin{extensionbox}
\textbf{Undo a sequence --- socks then shoes.} Two elimination steps clear the first column of a
$3\times3$: first $E_{21}=\begin{bsmallmatrix} 1&0&0\\-2&1&0\\0&0&1 \end{bsmallmatrix}$, then
$E_{31}=\begin{bsmallmatrix} 1&0&0\\0&1&0\\-5&0&1 \end{bsmallmatrix}$, so the combined step is
$E_{31}E_{21}$. To undo them you reverse the order:
$(E_{31}E_{21})^{-1}=E_{21}^{-1}E_{31}^{-1}$ (take shoes off before socks). Write out that product.
What do you notice about where $2$ and $5$ land?
\ansline{$E_{21}^{-1}E_{31}^{-1}=\begin{bsmallmatrix} 1&0&0\\2&1&0\\0&0&1 \end{bsmallmatrix}\begin{bsmallmatrix} 1&0&0\\0&1&0\\5&0&1 \end{bsmallmatrix}=\begin{bsmallmatrix} 1&0&0\\2&1&0\\5&0&1 \end{bsmallmatrix}$.
The multipliers $2$ and $5$ slot straight into positions $(2,1)$ and $(3,1)$ --- no interference. That
lower-triangular matrix of multipliers is the $L$ of $A=LU$, coming in 2.3.}
\end{extensionbox}

\begin{spiralbox}
\textbf{Coming next --- Lesson 2.3: Matrix Computations and $A=LU$.} You just built the undo-matrices
$E^{-1}$ that \emph{add} the multipliers back. Collect them into one lower-triangular matrix $L$ and
you get the clean factorization $A=LU$ --- elimination, packaged.
\end{spiralbox}

\end{document}
